\documentclass[10pt,a4paper]{article}
\usepackage[margin=1.7cm,top=1.5cm,bottom=1.6cm]{geometry}
\usepackage[T1]{fontenc}
\usepackage{lmodern}
\usepackage{microtype}
\usepackage{booktabs}
\usepackage{tabularx}
\usepackage{enumitem}
\usepackage{xcolor}
\usepackage{hyperref}
\hypersetup{colorlinks=true,linkcolor=black,urlcolor=blue!50!black,citecolor=black}
\setlist{itemsep=1pt,topsep=2pt,parsep=0pt,leftmargin=1.2em}
\setlength{\parskip}{3pt}
\setlength{\parindent}{0pt}
\newcommand{\M}{\,M}
\renewcommand{\arraystretch}{1.05}

\title{\vspace{-0.6cm}\textbf{Post-Quantum DTLS 1.3 with Mutual Authentication\\on a Bare-Metal RISC-V IoT Device}\\[2pt]
\large Inter IIT Tech Meet 14.0 --- QTrino Labs problem statement --- Technical Report}
\author{Nikhil Pothuganti, IIT Kanpur}
\date{\vspace{-0.9cm}}

\begin{document}
\maketitle

\section{Overview}
We built a DTLS~1.3 client for the provided LiteX/VexRiscv SoC (RV32IM, 4\,KB direct-mapped I- and
D-caches, simulated cycle-accurately with Verilator) and a host server, both on wolfSSL~5.8.4.
The device runs bare metal with our own interrupt-driven Ethernet/ARP/IPv4/UDP layer on the LiteEth MAC.
\begin{itemize}
  \item \textbf{Key exchange:} ML-KEM-512 (FIPS~203). It is the only group offered, with no classical fallback.
  \item \textbf{Authentication:} ML-DSA-44 (FIPS~204) certificates on \emph{both} sides. The server requires a client
    certificate, the device verifies the server's chain and hostname, and certificate validity dates are checked.
  \item \textbf{Records:} \texttt{TLS\_CHACHA20\_POLY1305\_SHA256}, with a HelloRetryRequest cookie against DoS amplification.
\end{itemize}
All cycle counts are read from the RISC-V \texttt{cycle} CSR. Times in ms assume a 100\,MHz clock.

\begin{table}[h]
\centering\small
\begin{tabular}{@{}lrrrr@{}}
\toprule
Metric (25 consecutive sessions, all successful) & Median & p90 & Min & Max \\
\midrule
DTLS 1.3 handshake (\texttt{wolfSSL\_connect}) & \textbf{12.44\M cyc = 124\,ms} & 15.9\M & 10.4\M & 35.7\M$^{a}$ \\
Full session establishment (context setup + handshake) & \textbf{13.89\M cyc = 139\,ms} & 17.4\M & 11.9\M & 37.2\M \\
\quad ML-DSA-44 sign (device CertificateVerify) & 4.43\M & 7.75\M & 2.80\M & 27.4\M \\
\quad ML-DSA-44 verify (server CertificateVerify) & 2.52\M & & & \\
\quad ML-KEM-512 keygen + decapsulation & 0.57\M\ + 0.88\M & & & \\
\quad Waiting on the network (inside the handshake) & 0.10\M & & & \\
Post-handshake throughput (48 records $\times$ 1400\,B) & \textbf{86 cyc/B = 1.14\,MB/s} & & & \\
ROM (\texttt{boot.bin}) / static RAM (.data+.bss) & \textbf{271.5\,KB} / 30.6\,KB & & & \\
Peak heap / peak stack & \textbf{98.5\,KB}$^{b}$ / 2.1\,KB & & & \\
\bottomrule
\end{tabular}

\vspace{2pt}
\footnotesize $^{a}$ML-DSA signing repeats a random number of times (rejection sampling; 4.25 tries on average for
ML-DSA-44). The maximum needed about 30 tries. Excluding the signature, the handshake is a stable ${\approx}8$\M\ cycles.
$^{b}$A \texttt{LOWMEM=1} build signs with wolfSSL's small-memory ML-DSA (61\,KB heap) at about $2\times$ the average signing time.
\caption{Final measurements (\texttt{scripts/measure.sh}, \texttt{results/measure\_20261116\_155400/}).}
\end{table}

\section{Architecture and integration}
\begin{itemize}
\item \textbf{Firmware.} \texttt{dtls\_client.c} drives wolfSSL through user I/O callbacks over our UDP layer
  (\texttt{net.c}).
  \begin{itemize}
    \item The Ethernet ISR copies each frame out of the MAC's two hardware RX slots into a 16-frame ring and answers
      ARP itself. Without that, the host's neighbour entry expired during long PQC computations and the kernel held
      back our packets.
    \item Time comes from the cycle counter. Entropy is a hook (\texttt{wolf\_port.c}) feeding wolfCrypt's Hash-DRBG
      (see Limitations).
    \item The device credentials are a CA certificate, the server certificate, and the device certificate and key.
      They are embedded as DER, and the key is in expanded-only form, so no key generation runs at load.
  \end{itemize}
\item \textbf{Configuration.} wolfSSL is set up entirely in \texttt{user\_settings.h}, with every option justified
  in comments. Changes outside wolfSSL's configuration: two small, documented patches to \texttt{dilithium.c}
  (Sec.~\ref{sec:opt}) and one \texttt{INCLUDE} line in the kit's linker script.
\item \textbf{Server.} The host server (native wolfSSL) requires and verifies the client certificate.
  It exports TLS secrets, so the capture decrypts completely in Wireshark/tshark~$\geq$4.4:
  CertificateRequest, both Certificate and CertificateVerify messages, Finished, ACK and the application data
  are all visible (\texttt{results/capture\_final/}).
\end{itemize}

\textbf{Bring-up issues and their root causes.} Each of these is useful to know for this platform.
\begin{enumerate}[label=(\arabic*)]
  \item VexRiscv traps on misaligned accesses and LiteX's handler loops silently. We found this with our own trap
    catcher; the fix is \texttt{WOLFSSL\_USE\_ALIGN}.
  \item Under DTLS~1.3 on a slow client, stale ACKs triggered endless fast retransmits. The fix is a larger
    record-number history and a 4\,s initial timer, as RFC~7925 recommends for constrained nodes.
  \item wolfSSL~5.8.4 bug: compiling out ML-KEM-1024 breaks ML-KEM-512 group support (a misplaced \texttt{break}).
  \item \texttt{dtls13.c} does not build without AES.
  \item AES on RISC-V defaults to cache-line ``touching'', which took 79\% of the cycles.
\end{enumerate}

\section{Optimisation}\label{sec:opt}
Every step was measured with the same harness. One clean build plus one session is logged to
\texttt{results/bench\_log.txt}, per-primitive microbenchmarks have known-answer tests (\texttt{BENCH=1}), and
PC-sampling profiles come from a timer interrupt.

The key platform fact: with \textbf{4\,KB direct-mapped caches}, code size \emph{is} speed.
\begin{itemize}
  \item \texttt{-O2} made the handshake slower: setup went from 5.7 to 14.0\M\ cycles, because the I-cache thrashed.
  \item Several later wins came from memory layout rather than instruction count.
\end{itemize}

\begin{table}[h]
\centering\small
\begin{tabularx}{\textwidth}{@{}>{\raggedright\arraybackslash}Xrrr@{}}
\toprule
Step (cumulative) & Handshake & Bulk & ROM \\
\midrule
First working build: AES-256-GCM-SHA384, \texttt{-Os} & 58--70\M & 744 cyc/B & 424\,KB \\
AES-128-GCM-SHA256 (SHA-384 is 64-bit arithmetic on RV32; AES-128 has 10 rounds) & 40.0\M & 584 & 403 \\
4-bit GHASH table instead of bit-serial GF($2^{128}$) multiply & 36.3\M & 361 & 404 \\
Network: no D-cache flush in ISR (MAC buffers are uncached I/O), 50\,ms gap ACK (RFC~9147 \S7.1), no IPv6 on tap & 29.0\M & 365 & 405 \\
Expanded-only ML-DSA key (OpenSSL's seed+key PKCS\#8 made wolfSSL re-run keygen): setup 5.8$\to$2.3\M & 26.5\M & 356 & 404 \\
ChaCha20-Poly1305: 87 vs.\ 320 cyc/B for AES-GCM (T-tables fill the whole D-cache); also constant-time & 25.6\M & 117 & 405 \\
Word-wise \texttt{memcpy}/\texttt{memset} via linker \texttt{--wrap} (picolibc is built \texttt{-Os}: byte loops) & 24.6\M & 89 & 406 \\
Keccak-f[1600] as explicit 32-bit lane halves, one fused round (23.1k$\to$16.2k cyc) & $\approx$20\M$^{c}$ & 87 & 403 \\
\texttt{-ffunction-sections}+GC, release config, one parameter set, no session cache (BSS 81$\to$26\,KB) & & 86 & \textbf{260} \\
Pinned server certificate: skips the CA-signature verify (CertificateVerify is still checked) & & & 267 \\
PQC-tuned layout: padded ML-DSA scratch buffers (they aliased in the D-cache; sign $-8\%$) & & & \\
PQRV RV32IM assembly~[1]: Keccak (15.2k cyc) and ML-DSA NTT (15.1k vs.\ 35k); signatures byte-identical & \textbf{12.4\M}$^{d}$ & \textbf{86} & 271 \\
\bottomrule
\end{tabularx}

\vspace{2pt}
\footnotesize $^{c}$From this point single-run totals are dominated by the random number of ML-DSA signing tries. Changes were compared with per-operation timers and fixed-seed microbenchmarks.
$^{d}$Median of 25 sessions.
\caption{Optimisation steps and their measured effect (cycles, single sessions unless noted).}
\end{table}

\textbf{Things we tried and rejected (all measured):}
\begin{itemize}
  \item Global or per-file \texttt{-O2} on SHA-3, ML-DSA or SHA-256: all slower because of the I-cache.
  \item AES small tables: 1.8$\times$ slower.
  \item wolfSSL's SHA-3 variants.
  \item Signed-Montgomery rewrite of 241 call sites: +15\%, a code-layout effect.
  \item A custom SHA-256: only $-4\%$, because it is instruction-bound and the handshake must hash ${\approx}26$\,KB.
  \item \texttt{ASSIGN\_KEY}: it breaks verification.
  \item ML-DSA's \texttt{SIGN\_CHECK\_Y/W0}: not FIPS~204-conformant.
\end{itemize}

\textbf{Memory.} The heap peak is inside ML-DSA signing: 44\,KB of TLS state plus 51\,KB for the signer's
matrix $A$ and vectors. wolfSSL's low-memory signer brings the peak to 61\,KB at a $2\times$ average signing cost,
so we ship it as a build option.

\textbf{Where the cycles go.} This is a median session (12.44\M\ cycles). The PQC operations come from
per-operation timers, SHA-256 from PC sampling.

\begin{center}\small
\begin{tabular}{@{}lrr@{}}
\toprule
Component & Cycles & Share \\
\midrule
ML-DSA-44 sign (device CertificateVerify; random number of tries) & 4.43\M & 36\% \\
ML-DSA-44 verify (server CertificateVerify) & 2.52\M & 20\% \\
SHA-256: transcript over ${\approx}$17\,KB of PQC handshake messages, plus ${\approx}$30 HKDF/HMAC derivations & ${\approx}$2.2\M & 18\% \\
ML-KEM-512 decapsulation (keygen, 0.57\M, is done during context setup) & 0.88\M & 7\% \\
ChaCha20-Poly1305 records, X.509/ASN.1 parsing, DTLS state machine, network wait & ${\approx}$2.4\M & 19\% \\
\bottomrule
\end{tabular}
\end{center}

Keccak (SHAKE) is about a third of all of it: ${\approx}$320 permutations per handshake, needed to expand the
ML-KEM and ML-DSA matrices and to sample ML-DSA's masks. This is why the Keccak work paid off across the board.

\section{Network and test methodology}
In the simulator, the PHY delivers frames at one byte per CPU cycle, back to back. The MAC has only two RX slots,
so the third frame of a server burst was lost while the ISR copied the first.
\begin{itemize}
  \item \textbf{Link model.} We shape \texttt{tap0} with \texttt{netem} at 2\,Mbit/s. The simulator executes
    ${\approx}$2\,M cycles per wall-clock second, so this equals a 100\,Mbit/s link at 100\,MHz \emph{in device cycles}.
  \item \textbf{Loss recovery.} Remaining losses of small back-to-back frames are recovered by the gap ACK; the
    decrypted capture shows one such recovery.
  \item \textbf{Server timer.} The server's initial retransmission timer is 10\,s, because it runs in real time
    while the simulated device is ${\approx}50\times$ slower than a 100\,MHz part.
  \item \textbf{Negative tests.} \texttt{scripts/verify.sh} checks that the device rejects a server certificate from
    a rogue CA ($-188$), one with the wrong hostname ($-322$) or an expired one ($-151$). It also checks that the server
    rejects a device certificate from a rogue CA ($-188$) or a missing one ($-345$). All rejections are clean on both
    ends, with no hang.
\end{itemize}

\section{Security discussion and limitations}
\begin{itemize}
  \item \textbf{Pinning.} The device trusts the CA \emph{and} pins the server certificate. If the server presents
    exactly that certificate, the CA signature on it is not re-verified (one ML-DSA verify, ${\approx}$2.5\M\ cycles).
    Possession of the key is still proven by CertificateVerify, and any other certificate is fully validated
    (RFC~7925 IoT profile). It can be disabled with \texttt{NOPIN=1}.
  \item \textbf{Cookie.} We keep the HRR cookie. Removing it would save one round trip and ${\approx}$0.3\M\ cycles,
    but the server would then do PQC work for spoofed ClientHellos.
  \item \textbf{Entropy (main open item).} The SoC has no TRNG. Our seed hook mixes the cycle counter into
    wolfCrypt's Hash-DRBG, which is \emph{not} secure entropy. A real deployment needs a hardware TRNG (e.g.\ a
    ring-oscillator source) feeding the same hook.
  \item \textbf{Time.} There is no RTC; the clock starts at the firmware build time for certificate validity checks.
  \item \textbf{Assembly code.} The PQRV assembly uses \texttt{gp}/\texttt{tp} as scratch registers. The build fails
    if the image ever defines a global pointer.
\end{itemize}

\vspace{2pt}
{\small [1] J.~Zhang, Y.~Yan, J.~Huang, Ç.\,K.~Koç, ``Optimized Software Implementation of Keccak, Kyber, and
Dilithium on RV\{32,64\}IM\{B\}\{V\}'', \emph{IACR TCHES} 2025(1), pp.~632--655. Code (MIT):
\url{https://github.com/Ji-Peng/PQRV}.}

\end{document}
